{\color{red}\noindent[\textbf{Fix P4} (fix\_poles\_dglap Steps 8--11): \textbf{explicit form} of the same term, with the $\xi\to0$ subtraction written out. Notation: $c_{0}\equiv2e^{-\gamma_{E}}$, $\mathcal{C}\equiv\frac{1}{(2\pi)^{3}}\frac{g^{2}}{2\pi^{2}}\frac{g^{2}}{8\pi^{2}}\frac{1}{k^{+}}=\mathcal{N}_{0}\frac{\alpha_{s}}{2\pi}$, $\mathcal{N}_{0}\equiv\frac{1}{(2\pi)^{3}}\frac{g^{2}}{2\pi^{2}}\frac{1}{k^{+}}$ (LO prefactor), $\beta_{0}=\frac{11N_{c}}{3}$, and the LO structure
$\mathcal{O}(y,y')\equiv\frac{(x-y)\cdot(x'-y')}{(x-y)^{2}(x'-y')^{2}}\big[U^{ce'}(y')-U^{ce'}(x')\big]\big[U^{ce}(y)-U^{ce}(x)\big]\rho^{e'}(x')\rho^{e}(x)$, $\widehat{L}(q)\equiv\int_{x,y,x',y'}e^{-iq\cdot(y-y')}\mathcal{O}(y,y')$, $\frac{d^{3}N_{\rm LO}}{d^{3}k}=\mathcal{N}_{0}\widehat{L}(k)$, $\Pi(\xi)\equiv\xi^{-2}\widehat{L}(k/\xi)$.]
\begin{align}
\frac{d^{3}N_{\rm DGLAP}}{d^{3}k}&=-\mathcal{C}\int_{0}^{1}d\xi\,\Big\{\mathcal{P}_{gg}(\xi)\,\mathcal{F}(\xi)-2N_{c}\,\frac{1-\xi}{\xi}\,\mathcal{F}(0)\Big\}\,,\qquad
\mathcal{F}(\xi)\equiv\frac{1}{\xi^{2}}\int_{x,y,x',y'}e^{-ik\cdot(y-y')\frac{1}{\xi}}\,\log\frac{(y'-y)^{2}\mu^{2}_{\overline{MS}}}{c_{0}^{2}\,\xi^{2}}\;\mathcal{O}(y,y')\,,\notag\\
\mathcal{P}_{gg}(\xi)&=2N_{c}\Big[\frac{\xi}{(1-\xi)_{+}}+\frac{1-\xi}{\xi}+\xi(1-\xi)\Big]+\frac{\beta_{0}}{2}\,\delta(1-\xi)\,,\qquad\int_{0}^{1}d\xi\,\frac{\xi\,f(\xi)}{(1-\xi)_{+}}\equiv\int_{0}^{1}d\xi\,\frac{\xi f(\xi)-f(1)}{1-\xi}\,,\label{eq:fx-DGLAP}\\
\mathcal{F}(0)&=-\frac{4\pi}{k^{2}}\,\mathcal{X}\,,\qquad\mathcal{X}\equiv\int_{x,x',z}\mathcal{O}(z,z)=\int_{x,x',z}\frac{(x-z)\cdot(x'-z)}{(x-z)^{2}(x'-z)^{2}}\big[U^{ce'}(z)-U^{ce'}(x')\big]\big[U^{ce}(z)-U^{ce}(x)\big]\rho^{e'}(x')\rho^{e}(x)\,.\notag
\end{align}
\noindent[\textbf{Calculation.} $\log\frac{r^{2}\mu^{2}}{c_{0}^{2}\xi^{2}}=2\gamma_{E}+\log\frac{r^{2}\mu^{2}}{4\xi^{2}}$, so $\mathcal{F}(\xi)$ is the author's function under $\int\frac{d\xi}{\xi^{2}}$; $\widetilde{\mathcal{P}}_{gg}$ with $\int_{0}^{1}d\xi\big[\frac{1-\xi}{\xi}\big]_{+}\mathcal{F}\equiv\int_{0}^{1}d\xi\,\frac{1-\xi}{\xi}\big[\mathcal{F}(\xi)-\mathcal{F}(0)\big]$ is $\mathcal{P}_{gg}$ with the subtraction $-2N_{c}\frac{1-\xi}{\xi}\mathcal{F}(0)$: Eq.~\eqref{eq:fx-DGLAP} is the author's term, unchanged in content. $\mathcal{F}(0)$: with $y'-y=\xi s$, $\mathcal{F}(\xi)\to\int d^{2}s\,e^{ik\cdot s}\log s^{2}\int_{x,x',z}\mathcal{O}(z,z)$ and $\int d^{2}s\,e^{-ik\cdot s}\log s^{2}=-\frac{4\pi}{k^{2}}$ ($k\neq0$). The $\xi\to0$ (BFKL) subtraction carries no pole ($\Pi(\xi)\to0$ for $k\neq0$).]

\noindent[\textbf{NOTE (subtraction term).} fix\_poles\_dglap Step 11 writes the subtraction as $-\frac{2N_{c}}{\xi}\mathcal{F}(0)$ instead of $-2N_{c}\frac{1-\xi}{\xi}\mathcal{F}(0)$. The two differ by $\int_{0}^{1}d\xi\big[\frac{2N_{c}}{\xi}-2N_{c}\frac{1-\xi}{\xi}\big]\mathcal{F}(0)=2N_{c}\mathcal{F}(0)$:
\begin{align}
-\mathcal{C}\int_{0}^{1}d\xi\Big\{\mathcal{P}_{gg}\mathcal{F}-\frac{2N_{c}}{\xi}\mathcal{F}(0)\Big\}=\text{Eq.~\eqref{eq:fx-DGLAP}}\;\underbrace{-\,2N_{c}\,\mathcal{C}\,\frac{4\pi}{k^{2}}\,\mathcal{X}}_{=\;-\frac{1}{(2\pi)^{3}}\frac{g^{4}N_{c}}{2\pi^{3}}\frac{1}{k^{2}}\frac{1}{k^{+}}\mathcal{X}}\,,\notag
\end{align}
and the underbraced term is exactly the remnant ($-\frac{1}{k^{+}}\times$BFKL structure) of Group II Row IV (its second term). The Step 11 form is therefore DGLAP $+$ remnant (the KLSZ ``double plus'' form, \texttt{lit\_dglap\_rc}); Step 11 also says that the remnant stays among the finite terms of Row IV, which would count it twice. Here the author's convention is kept: subtraction $2N_{c}\frac{1-\xi}{\xi}\mathcal{F}(0)$ inside DGLAP, remnant in Group II Row IV. \emph{Question for the author:} which of the two equivalent writings is wanted in the paper (if the $\frac{2N_{c}}{\xi}$ form, the remnant has to be deleted from Row IV).]

\noindent[\textbf{Scheme} (fix\_poles\_dglap Steps 8--10; reviewer's point P4). (a) In full dimensional regularisation ($d=2-\epsilon$ also for the separation $u=w'-w$ of the pair, rescaled as $y'-y=\xi u$ in $d$ dimensions: $\int d^{d}u\,(u^{2})^{\epsilon/2}e^{iku}f(\xi u)=\xi^{-2}\int d^{d}s\,(s^{2})^{\epsilon/2}e^{i\frac{k}{\xi}s}f(s)$) the \emph{bare} collinear term is $\xi^{-2}\big[\frac{2}{\epsilon}+\log\frac{(y'-y)^{2}\mu^{2}}{c_{0}^{2}}\big]_{d}\times$(LO structure of the parent): in the bare term there is no $\log\xi^{2}$ (the reviewer is right on this point; the author's $\log\frac{(y'-y)^{2}\mu^{2}}{4\xi^{2}}$ arises from continuing only $r=w-z$). (b) What survives the subtraction depends on the $\xi$ power of the counterterm. For an $\overline{\rm MS}$ fragmentation function defined as a \emph{number density} with $d$-dimensional collinear kinematics ($P=z\,p$ in all $d$ components, $\frac{dN_{h}}{dP^{+}d^{d}P}=\int\frac{dz}{z^{1+d}}D(z)\frac{dN_{g}}{dk^{+}d^{d}k}\big|_{k=P/z}$, pure-pole renormalisation $D_{0}=[\delta(1-\xi)+\frac{\alpha_{s}}{2\pi}\frac{2}{\epsilon}\mathcal{P}_{gg}+\dots]\otimes D(\mu)$), the partonic counterterm is $+\mathcal{C}\frac{2}{\epsilon}\int d\xi\,\mathcal{P}_{gg}(\xi)\,\xi^{-d}\,\widehat{L}_{d}(k/\xi)$; $\xi^{-d}=\xi^{-2}(1+\epsilon\log\xi)$ restores the $\log\xi^{2}$: bare $-$ CT $=$ the author's form, Eq.~\eqref{eq:fx-DGLAP} [fix\_poles\_dglap \texttt{s3}: $-4.38223062$ against $-4.38221572$ in $d=2\pm0.002$]. With the two-dimensional Jacobian $\xi^{-2}$ (Chirilli--Xiao--Yuan) one gets instead $\log\frac{(y'-y)^{2}\mu^{2}}{c_{0}^{2}}$ (CXY form). \textbf{Definition} (\textbf{author's answer (date 2026-11-20):} ``Yes'', the fragmentation functions are $\overline{\rm MS}$ number densities, i.e.\ the author's $\log\frac{(y'-y)^{2}\mu^{2}}{4\xi^{2}}$ form): the fragmentation functions are $\overline{\rm MS}$ number densities, renormalised with a pure pole and the $d$-dimensional collinear Jacobian $\xi^{-d}$ (this definition satisfies the momentum sum rule); Eq.~\eqref{eq:fx-DGLAP} is the collinear term for this definition. (c) \emph{Remark} (the CXY convention, given only so that it can be recovered). The difference to the CXY form is
\begin{align}
\Delta\equiv\frac{d^{3}N_{\rm DGLAP}}{d^{3}k}-\frac{d^{3}N_{\rm DGLAP}^{\rm CXY}}{d^{3}k}=+\,\mathcal{C}\int_{0}^{1}d\xi\,\mathcal{P}_{gg}(\xi)\,\log\xi^{2}\;\Pi(\xi)\,,\qquad \Pi(\xi)=\frac{1}{\xi^{2}}\int_{x,y,x',y'}e^{-ik\cdot(y-y')\frac{1}{\xi}}\,\mathcal{O}(y,y')\,,\label{eq:fx-Delta}
\end{align}
where $N^{\rm CXY}_{\rm DGLAP}$ is Eq.~\eqref{eq:fx-DGLAP} with $\log\frac{(y'-y)^{2}\mu^{2}}{c_{0}^{2}\xi^{2}}\to\log\frac{(y'-y)^{2}\mu^{2}}{c_{0}^{2}}$ ($\mathcal{F}(0)$ is the same in both forms). $\Delta$ has no $\delta(1-\xi)$ part, is integrable at $\xi\to0$, does not depend on $\mu$; equivalently $D^{\rm CXY}=\big[\delta(1-\xi)+\frac{\alpha_{s}}{2\pi}\mathcal{P}_{gg}(\xi)\log\xi^{2}\big]\otimes D$. Since $\int_{0}^{1}d\xi\,\xi\,\mathcal{P}_{gg}\log\xi^{2}=4N_{c}\big[\frac54-\zeta_{2}-\frac34-\frac{7}{144}\big]=-14.3225$ ($N_{c}=3$), the two definitions cannot both satisfy the momentum sum rule. (d) $\mu$ dependence: $\partial_{\log\mu^{2}}N_{\rm DGLAP}=-\mathcal{C}\int_{0}^{1}d\xi\,\mathcal{P}_{gg}\,\Pi$, the DGLAP evolution of $D_{g}(\mu)$. [\textbf{Author's answer (date 2026-11-20):} the number-density definition is confirmed; the former NOTE on the choice of the FF definition is answered and removed. Not done (a cross-check only, not a choice of convention): an independent comparison with a standard NLO collinear calculation of $pp\to gX$.]]

\noindent[\textbf{Poles} (fix P3; fix\_poles\_dglap Step 4). The collinear pole of Group II Row IV (large daughter separation, IR) multiplies the real splitting function $\mathcal{P}_{gg+}=\mathcal{P}_{gg}-\frac{\beta_{0}}{2}\delta(1-\xi)$ (no $\delta(1-\xi)$: a real emission): $-\mathcal{C}\frac{2}{\epsilon}\int d\xi\,\mathcal{P}_{gg+}\Pi=-\mathcal{C}\frac{2}{\epsilon}[\mathcal{P}_{gg}\otimes\Pi]+\frac{\alpha_{s}\beta_{0}}{4\pi}\frac{2}{\epsilon}\,\frac{d^{3}N_{\rm LO}}{d^{3}k}$. The first part is absorbed by the fragmentation counterterm with the full $\mathcal{P}_{gg}$; the second is the pole of the ``RC'' term below and is collinear (IR). The $\xi\to0$ subtraction has no pole. Full table: \texttt{Divergent\_terms\_fixed}.]

\sloppy\noindent[\textbf{Both forms} (\textbf{author's answer (date 2026-11-20):} ``write it in both forms''; fix\_poles\_dglap Steps 12--14, \texttt{notes/fix\_tail.tex}). The DGLAP and RC terms are given in two exactly equivalent forms: the \emph{coordinate form}, Eq.~\eqref{eq:fx-DGLAP} and $N_{\rm RC}=+\mathcal{C}\frac{\beta_{0}}{2}\mathcal{F}(1)$ (RC term below), and the \emph{local (momentum-space) form}. With $\log\frac{r^{2}\mu^{2}}{c_{0}^{2}}=-\frac{1}{\pi}\int d^{2}l\,\frac{e^{-il\cdot r}-\theta(\mu-|l|)}{l^{2}}$ one has exactly
\begin{align}
\mathcal{F}(\xi)=-\log\frac{k^{2}}{\mu^{2}}\,\Pi(\xi)+\frac{1}{\xi^{2}}\,\mathcal{R}\Big(\frac{k}{\xi}\Big),\qquad
\mathcal{R}(q)\equiv-\frac{1}{\pi}\int d^{2}l\,\frac{\widehat{L}(q+l)-\theta(|q|-|l|)\,\widehat{L}(q)}{l^{2}},\qquad\lim_{\xi\to0}\xi^{-2}\mathcal{R}(k/\xi)=\mathcal{F}(0),\label{eq:fx-loc}
\end{align}
($\mathcal{R}$ does not depend on $\mu$). The local DGLAP and RC terms are
\begin{align}
\frac{d^{3}N^{\rm loc}_{\rm DGLAP}}{d^{3}k}\equiv\mathcal{C}\log\frac{k^{2}}{\mu^{2}}\int_{0}^{1}d\xi\,\mathcal{P}_{gg}(\xi)\,\Pi(\xi)\,,\qquad
\frac{d^{3}N^{\rm loc}_{\rm RC}}{d^{3}k}\equiv-\frac{\alpha_{s}\beta_{0}}{4\pi}\log\frac{k^{2}}{\mu^{2}}\,\mathcal{N}_{0}\widehat{L}(k)=-\mathcal{C}\,\frac{\beta_{0}}{2}\log\frac{k^{2}}{\mu^{2}}\,\Pi(1)\,,\notag
\end{align}
and, inserting Eq.~\eqref{eq:fx-loc} into the coordinate form (with the subtraction of Eq.~\eqref{eq:fx-DGLAP}),
\begin{align}
\frac{d^{3}N_{\rm DGLAP}}{d^{3}k}&=\frac{d^{3}N^{\rm loc}_{\rm DGLAP}}{d^{3}k}-\mathcal{C}\int_{0}^{1}d\xi\Big\{\mathcal{P}_{gg}(\xi)\,\frac{\mathcal{R}(k/\xi)}{\xi^{2}}-2N_{c}\frac{1-\xi}{\xi}\,\mathcal{F}(0)\Big\}\,,\qquad
\frac{d^{3}N_{\rm RC}}{d^{3}k}=\frac{d^{3}N^{\rm loc}_{\rm RC}}{d^{3}k}+\mathcal{C}\,\frac{\beta_{0}}{2}\,\mathcal{R}(k)\,,\label{eq:fx-locsplit}\\
\frac{d^{3}N_{\rm DGLAP}}{d^{3}k}+\frac{d^{3}N_{\rm RC}}{d^{3}k}&=\frac{d^{3}N^{\rm loc}_{\rm DGLAP}}{d^{3}k}+\frac{d^{3}N^{\rm loc}_{\rm RC}}{d^{3}k}+\frac{d^{3}N_{\mathcal{R}}}{d^{3}k}\,,\qquad
\frac{d^{3}N_{\mathcal{R}}}{d^{3}k}\equiv-\mathcal{C}\int_{0}^{1}d\xi\Big\{\mathcal{P}_{gg+}(\xi)\,\frac{\mathcal{R}(k/\xi)}{\xi^{2}}-2N_{c}\frac{1-\xi}{\xi}\,\mathcal{F}(0)\Big\}\,,\label{eq:fx-equiv}
\end{align}
$\mathcal{P}_{gg+}=\mathcal{P}_{gg}-\frac{\beta_{0}}{2}\delta(1-\xi)$: the $\delta(1-\xi)$ parts $\mp\mathcal{C}\frac{\beta_{0}}{2}\mathcal{R}(k)$ of the two non-local pieces cancel. The local terms carry all the $\mu$ dependence and need no subtraction at $\xi=0$ ($\Pi(\xi)\to0$). The non-local remainder $N_{\mathcal{R}}$ is $\mu$-independent and belongs to the finite terms of Group II Row IV. [fix\_poles\_dglap Step 12 writes $N_{\mathcal{R}}$ with $\frac{2N_{c}}{\xi}\mathcal{F}(0)$; that is the $N_{\mathcal{R}}$ above plus the remnant of Group II Row IV (its second term), see the NOTE (subtraction term) above.]]

\noindent[\textbf{Large-distance behaviour of the two forms} (\textbf{author's answer (date 2026-11-20):} both forms; settled in fix\_tail). At large parent position $Y=\frac{y+y'}{2}$ the LO structure tends to $G_{\infty}(Y)/Y^{2}$ independently of $y'-y$, the logarithm survives the Fourier transform, and $\mathcal{F}(\xi)\to-\frac{4\pi}{k^{2}}G_{\infty}/Y^{2}$ for every $\xi$. \emph{Coordinate form:} the DGLAP and RC terms are each logarithmically IR divergent at large $Y$ (soft-parent region; tail coefficients in units $\mathcal{C}\frac{4\pi}{k^{2}}\int\frac{d^{2}Y}{Y^{2}}G_{\infty}$: DGLAP $+\frac{N_{c}}{6}$, RC $-\frac{11N_{c}}{6}$; in Group II Row IV also the remnant $-2N_{c}$ and the traceless tensor term $+\frac{N_{c}}{3}\cos2\varphi$, total $2N_{c}[-\frac{11}{6}+\frac16\cos2\varphi]$). Their sum with Group II Row I (second part) and Group III Row III is not divergent: the tail cancels term by term (\texttt{notes/fix\_tail.tex}: in units $2N_{c}\mathcal{C}\frac{4\pi}{k^{2}}\int\frac{d^{2}Y}{Y^{2}}G_{\infty}$, Group II Row IV $-\frac{11}{6}+\frac16\cos2\varphi$, Group II Row I second part $-\frac{11}{6}+\frac16\cos2\varphi$, Group III Row III second part with c.c.\ $+\frac{11}{3}-\frac13\cos2\varphi$; sum $0$, also in the BFKL and JIMWLK sectors), so the three rows must be integrated together over the pair position. \emph{Local form:} $N^{\rm loc}_{\rm DGLAP}$ and $N^{\rm loc}_{\rm RC}$ have no tail; the tails of the coordinate-form DGLAP and RC terms ($+\frac{N_{c}}{6}-\frac{11N_{c}}{6}$) are carried by the remainder $N_{\mathcal{R}}$ of Eq.~\eqref{eq:fx-equiv}, which (with the remnant and the tensor term of Group II Row IV, whose tails are the same in both forms) has to be integrated together with Group II Row I (second part) and Group III Row III.]}\par
